% 自写模板 — 华数杯全国大学生数学建模竞赛论文骨架 (精简版)
% 用法: xelatex main.tex (连编两遍); 或
%       python scripts/render_paper.py --competition huashubei --workspace <paper_workspace>
% 依据: 华数杯通行论文结构 (摘要页/正文/参考文献/附录)。
%       注意: 华数杯各届通知对封面字段、页边距、摘要字数等细节略有调整,
%       投稿前务必按当年竞赛通知核对并修改下方对应段落。
%       本文件为独立实现, 未参考任何第三方模板源码。
% 字体: 正文宋体小四, 章节标题黑体; xelatex 下 ctex 自动选用系统字体。

\documentclass[zihao=-4,a4paper]{ctexart}

% ---------- 页面 (默认 2.5cm; 若当年通知另有规定请修改) ----------
\usepackage{geometry}
\geometry{a4paper, top=2.5cm, bottom=2.5cm, left=2.5cm, right=2.5cm}

% ---------- 数学 ----------
\usepackage{amsmath,amssymb}

% ---------- 图表 ----------
\usepackage{graphicx}
\usepackage{float}
\usepackage{booktabs}
\usepackage{multirow}
\usepackage{longtable}
\usepackage{caption}

% ---------- 代码附录 ----------
\usepackage{listings}
\usepackage{xcolor}

% ---------- 其他 ----------
\usepackage{enumitem}
\usepackage[hidelinks]{hyperref}

% 章节标题: 黑体
\ctexset{
  section = {format = {\Large\bfseries\heiti}},
  subsection = {format = {\large\bfseries\heiti}},
  subsubsection = {format = {\bfseries\heiti}},
}

% 图表标题
\captionsetup{font={small,bf}, labelsep=quad}

% 代码样式
\lstset{
  basicstyle=\small\ttfamily,
  breaklines=true,
  frame=single,
  numbers=left,
  numberstyle=\tiny\color{gray},
  keywordstyle=\color{blue},
  commentstyle=\color{green!50!black},
  stringstyle=\color{red!60!black},
  showstringspaces=false,
  columns=flexible,
}

\begin{document}

% ============================================================
% 首页: 题目信息区 + 摘要 (华数杯通常不设独立封面;
% 若当年通知要求独立封面/承诺书, 在此前自行插入 titlepage)
% ============================================================
\pagenumbering{arabic}

\begin{center}
  {\Large\heiti 论文标题: [论文标题]}\\[0.3em]
  {\large 题号: [题号]\qquad 参赛队号: [参赛队号]\qquad 学校: [学校名称]}
\end{center}

\vspace{0.5em}

\begin{center}
  {\Large\heiti 摘\quad 要}
\end{center}

% ---- 摘要正文 (篇幅按当年通知, 一般 1 页以内) ----
% render_paper.py 注入锚: 下一行在渲染时被替换为 abstract 节内容; 手写论文时删除锚点行直接写正文
% __ABSTRACT__

\vspace{1em}
\noindent{\heiti 关键词:} [关键词1]; [关键词2]; [关键词3]

\newpage

% ============================================================
% 正文
% render_paper.py 注入锚: 下方成对 SECTIONS 锚行之间的
% 演示正文会被 sections/*.tex 自动替换; 手写论文时删除两个锚点行直接写
% ============================================================
% __SECTIONS__

\section{问题重述}

[两到三段交代问题背景与各子问任务, 不照抄题面。]

\section{问题分析}

[逐问分析求解思路, 给出总体技术路线。]

\section{模型假设与符号说明}

\begin{enumerate}
  \item [假设 1 及其合理性依据];
  \item [假设 2 及其合理性依据]。
\end{enumerate}

\begin{table}[H]
  \centering
  \caption{符号说明表}
  \begin{tabular}{cll}
    \toprule
    符号 & 含义 & 单位 \\
    \midrule
    $i$   & 决策变量编号 & --- \\
    $x_i$ & 第 $i$ 个决策变量 & [单位] \\
    $\lambda$ & 模型参数 & [单位] \\
    \bottomrule
  \end{tabular}
\end{table}

\section{模型建立与求解}

\subsection{[子问题 1] 的模型建立}

[模型推导过程, 示例公式如下。]

\begin{equation}
  \min \; Z = \sum_{i=1}^{n} c_i x_i
\end{equation}
\begin{equation}
  \text{s.t.}\quad \sum_{i=1}^{n} a_{ij} x_i \le b_j, \quad j = 1,2,\dots,m
\end{equation}

\subsection{模型求解与结果}

[求解算法与结果分析。]

\begin{table}[H]
  \centering
  \caption{[求解结果汇总]}
  \begin{tabular}{cccc}
    \toprule
    方案 & 指标 A & 指标 B & 目标值 \\
    \midrule
    方案 1 & [值] & [值] & [值] \\
    方案 2 & [值] & [值] & [值] \\
    方案 3 & [值] & [值] & [值] \\
    \bottomrule
  \end{tabular}
\end{table}

% 图占位示例 (实际使用时取消注释并替换路径)
\begin{figure}[H]
  \centering
  % \includegraphics[width=0.8\textwidth]{figures/result.png}
  \fbox{\parbox[c][6cm][c]{0.78\textwidth}{\centering [图占位: 技术路线图 / 结果可视化]}}
  \caption{[图题: 一句话说明该图支撑的结论]}
\end{figure}

\section{灵敏度分析}

[关键参数扰动与结论。]

\section{模型评价与推广}

[优点、缺点与改进方向。]

% __END_SECTIONS__

% ============================================================
% 参考文献 (GB/T 7714 顺序编码制, 手写条目示例)
% ============================================================
\begin{thebibliography}{9}
  \bibitem{ref1} 作者. 论文题目[J]. 期刊名, 2024, 20(3): 1-10.
  \bibitem{ref2} 作者. 书名[M]. 北京: 出版社, 2023.
  \bibitem{ref3} Author A. Paper title[J]. Journal Name, 2022, 15(2): 100-110.
\end{thebibliography}

% ============================================================
% 附录: 程序代码
% ============================================================
\appendix

\section{程序代码}

\begin{lstlisting}[language=Python, caption={求解脚本示例 solve.py}]
import numpy as np
from scipy.optimize import linprog

# model parameters
c = np.array([2.0, 3.0, 1.5])          # cost coefficients
A_ub = np.array([[1.0, 1.0, 0.0],
                 [0.0, 2.0, 1.0]])
b_ub = np.array([40.0, 50.0])

res = linprog(c, A_ub=A_ub, b_ub=b_ub,
              bounds=[(0, None)] * 3, method="highs")
print("optimal value:", res.fun)
print("solution:", res.x)
\end{lstlisting}

\end{document}
